\documentclass[14pt,a4paper,oneside]{extbook}
\usepackage[utf8]{vietnam}
\usepackage{amsmath,amssymb,mathrsfs,fancyhdr,enumerate,multirow,makecell}
\usepackage{tkz-euclide,tikz-3dplot,tikz,tkz-tab}
\usetikzlibrary{arrows.meta,calc,intersections,angles,quotes,backgrounds,shapes.geometric,patterns,through,shadows,positioning}
\usepackage{yhmath}
\usepackage[top=0.7cm, bottom=0.7cm, left=1.0cm, right=1.0cm, headheight=16pt, headsep=3mm, footskip=5mm]{geometry}
\usepackage[hidelinks,unicode]{hyperref}
\usepackage{qrcode}
\usepackage[many]{tcolorbox}
\usepackage{fontawesome}

\pagestyle{fancy}
\fancyhf{}
\fancyhead[L]{\footnotesize\bfseries\color{teal!85!black}\faBook~SỔ TAY KIẾN THỨC TOÁN 5}
\fancyhead[R]{\footnotesize\bfseries\color{blue!85!black}Lớp Toán Thầy Lĩnh -- loptoanthaylinh.vn}
\fancyfoot[C]{\footnotesize\bfseries\color{gray!80!black}-- Sổ tay kiến thức: Bài 27. Đường tròn. Chu vi và diện tích hình tròn --}
\renewcommand{\headrulewidth}{0.4pt}

\begin{document}

% --- BANNER TIÊU ĐỀ BÀI HỌC ---
\noindent
\begin{tcolorbox}[
  enhanced,
  colback=teal!8!white,
  colframe=teal!80!black,
  arc=3mm,
  boxrule=1.1pt,
  width=\linewidth,
  halign=center,
  top=2.2mm, bottom=2.2mm,
  left=3mm, right=3mm,
  drop fuzzy shadow
]
  {\large\bfseries\color{teal!95!black}\faBookmark\ BÀI 27. ĐƯỜNG TRÒN. CHU VI VÀ DIỆN TÍCH HÌNH TRÒN}\\[2pt]
  {\scriptsize\color{black!75}\itshape Bán kính r, đường kính d, số Pi (3,14), công thức Chu vi và Diện tích hình tròn trên 1 trang A4}
\end{tcolorbox}

\vspace{3.2mm}
% --- KHỐI 1: KHÁI NIỆM ĐƯỜNG TRÒN VÀ HÌNH TRÒN ---
\noindent
\begin{minipage}[t]{0.56\textwidth}
  \begin{tcolorbox}[
    enhanced,
    title={\footnotesize\bfseries\color{white}\faCircleO\ 1. Đường tròn và Hình tròn},
    colback=blue!2!white,
    colframe=blue!75!black,
    arc=2.5mm, boxrule=0.8pt, top=1.8mm, bottom=1.8mm, left=2.2mm, right=2.2mm
  ]
    \footnotesize
    $\bullet$ \textbf{Đường tròn tâm $O$ bán kính $r$:} Đường viền bao quanh.\\
    $\bullet$ \textbf{Hình tròn:} Gồm đường tròn và toàn bộ phần mặt phẳng nằm bên trong đường tròn đó.\\
    $\bullet$ \textbf{Bán kính ($r$):} Đoạn thẳng nối tâm $O$ với điểm trên đường tròn.\\
    $\bullet$ \textbf{Đường kính ($d$):} Đoạn thẳng đi qua tâm $O$ nối hai điểm trên đường tròn: \textbf{$d = 2 \times r$}.
  \end{tcolorbox}
\end{minipage}%
\hfill
\begin{minipage}[t]{0.42\textwidth}
  \begin{tcolorbox}[
    enhanced,
    title={\footnotesize\bfseries\color{white}\faEye\ Minh họa hình tròn},
    colback=teal!2!white,
    colframe=teal!75!black,
    arc=2.5mm, boxrule=0.8pt, top=1.8mm, bottom=1.8mm, left=1.5mm, right=1.5mm,
    halign=center
  ]
    \begin{tikzpicture}[scale=0.75]
      \draw[fill=teal!15, draw=teal!85!black, thick] (0,0) circle (1.2cm);
      \draw[blue!85!black, thick] (0,0) -- (1.2,0) node[midway, above, font=\tiny] {$r$};
      \draw[red!85!black, dashed, thick] (-1.2,0) -- (0,0);
      \filldraw (0,0) circle (1.5pt) node[below, font=\scriptsize] {$O$};
      \node[font=\tiny\bfseries, red!85!black] at (0,-0.7) {Đường kính $d = 2r$};
    \end{tikzpicture}\\[2pt]
    {\tiny\color{teal!95!black}\bfseries Số $\pi \approx 3,14$ dùng cho tiểu học}
  \end{tcolorbox}
\end{minipage}

\vspace{3mm}

% --- KHỐI 2: CÔNG THỨC CHU VI VÀ DIỆN TÍCH ---
\noindent
\begin{tcolorbox}[
  enhanced,
  title={\footnotesize\bfseries\color{white}\faCalculator\ 2. Bảng công thức Chu vi và Diện tích hình tròn},
  colback=purple!2!white,
  colframe=purple!70!black,
  arc=2.5mm, boxrule=0.8pt, top=1.8mm, bottom=1.8mm, left=2.5mm, right=2.5mm
]
  \footnotesize
  \begin{tabular}{p{0.48\textwidth}|p{0.48\textwidth}}
    \textbf{CHU VI HÌNH TRÒN ($C$)} & \textbf{DIỆN TÍCH HÌNH TRÒN ($S$)} \\ \hline
    $\bullet$ Theo đường kính: $\mathbf{C = d \times 3,14}$ & $\bullet$ Theo bán kính: $\mathbf{S = r \times r \times 3,14}$ \\
    $\bullet$ Theo bán kính: $\mathbf{C = r \times 2 \times 3,14}$ & \textit{(Bán kính nhân bán kính rồi nhân 3,14)} \\
    $\triangleright$ Tính ngược: $d = C : 3,14$ hoặc $r = C : 6,28$. & $\triangleright$ Tính ngược: $r \times r = S : 3,14 \rightarrow$ tìm $r$.
  \end{tabular}
\end{tcolorbox}

\vspace{3mm}

% --- KHỐI 3: SƠ ĐỒ TƯ DUY ---
\noindent
\begin{tcolorbox}[
  enhanced,
  title={\footnotesize\bfseries\color{white}\faSitemap\ Sơ đồ tư duy: Chu vi \& Diện tích hình tròn},
  colback=yellow!2!white,
  colframe=teal!85!black,
  arc=3mm, boxrule=0.9pt, top=1.8mm, bottom=1.8mm, left=1.5mm, right=1.5mm,
  halign=center
]
  \begin{tikzpicture}[
    scale=0.92,
    every node/.style={transform shape},
    root/.style={rectangle, rounded corners=5pt, fill=teal!85!black, text=white, font=\bfseries\footnotesize, inner sep=5pt, align=center},
    b1/.style={rectangle, rounded corners=4pt, fill=blue!15, draw=blue!70!black, text=blue!90!black, font=\bfseries\scriptsize, inner sep=4.5pt, align=center},
    b2/.style={rectangle, rounded corners=4pt, fill=orange!15, draw=orange!70!black, text=orange!90!black, font=\bfseries\scriptsize, inner sep=4.5pt, align=center}
  ]
    \node[root] (core) at (0, 0.75) {\faCircle\ HÌNH TRÒN};
    \node[b1] (n1) at (-4.7, -0.45) {\textbf{CHU VI: } $C = d \times 3,14 = 2 \times r \times 3,14$\\[2pt](Độ dài bánh xe lăn $1$ vòng)};
    \node[b2] (n2) at (4.7, -0.45) {\textbf{DIỆN TÍCH: } $S = r \times r \times 3,14$\\[2pt](Bề mặt bồn hoa, mặt đồng hồ)};

    \draw[teal!85!black, thick, -{Stealth[length=2mm]}] (core) -| (n1);
    \draw[teal!85!black, thick, -{Stealth[length=2mm]}] (core) -| (n2);
  \end{tikzpicture}
\end{tcolorbox}

\vspace{3mm}

% --- KHỐI 4: TOÁN BÁNH XE LĂN & BẪY TRẮC NGHIỆM ---
\noindent
\begin{minipage}[t]{0.54\textwidth}
  \begin{tcolorbox}[
    enhanced,
    title={\footnotesize\bfseries\color{white}\faBicycle\ Dạng toán bánh xe lăn trên đường},
    colback=green!3!white,
    colframe=green!60!black,
    arc=2.5mm, boxrule=0.8pt, top=1.8mm, bottom=1.8mm, left=2.2mm, right=2.2mm
  ]
    \footnotesize
    $\bullet$ Bánh xe lăn $1$ vòng đi được quãng đường bằng đúng \textbf{chu vi bánh xe}: $s_1 = C = d \times 3,14$.\\
    $\bullet$ Khi bánh xe lăn $n$ vòng: $\text{Quãng đường } S = C \times n$.\\
    $\bullet$ Số vòng lăn: $n = S : C$.
  \end{tcolorbox}
\end{minipage}%
\hfill
\begin{minipage}[t]{0.44\textwidth}
  \begin{tcolorbox}[
    enhanced,
    title={\footnotesize\bfseries\color{white}\faExclamationTriangle\ Mẹo tránh bẫy trắc nghiệm},
    colback=red!2!white,
    colframe=red!70!black,
    arc=2.5mm, boxrule=0.8pt, top=1.8mm, bottom=1.8mm, left=2.2mm, right=2.2mm
  ]
    \footnotesize
    $\bullet$ \textbf{Lỗi nhầm kinh điển}: Diện tích là $r \times r \times 3,14$, rất nhiều bạn nhầm thành $r \times 2 \times 3,14$ (đó là Chu vi!).\\
    $\bullet$ Cho đường kính $d$ mà tính diện tích: \textbf{Bắt buộc} phải tìm bán kính $r = d : 2$ trước.
  \end{tcolorbox}
\end{minipage}

\end{document}
